\chapter{Multiple Regression}\label{ch:b2:multiple-regression}

A regression with one regressor compares people who differ in that regressor and, usually, in much else. Adding
regressors compares people who differ in the first and are alike in the others. This chapter writes the model in matrix
form, derives the least-squares estimator, shows what ``holding the others constant'' means exactly, and measures what
leaving a variable out does to a coefficient. The chart shows the question on real data: the estimated return to a year
of schooling as the comparison is made more like for like.

\begin{openingchart}{Estimated effect of a year of schooling on log earnings in 1978, per cent, as controls are added one
  at a time, with 95 per cent confidence intervals; American men with positive earnings in 1975 and 1978 (13,055 men).}{CPS
  comparison sample of LaLonde (1986), as distributed by Dehejia and Wahba; our computation}
\begin{tikzpicture}
\begin{axis}[omrecord, width=0.72\linewidth, height=0.4\linewidth, xmin=-0.5, xmax=6.5, ymin=0, ymax=6.5,
    xtick={0,...,6}, xticklabels from table={figdata/bachelor-2/30-multiple-regression/out/steps.csv}{label},
    xticklabel style={rotate=30, anchor=north east, font=\scriptsize}, ylabel={Per cent per year}]
  \addplot[omDef, thick, mark=*, mark size=1.6pt, error bars/.cd, y dir=both, y explicit]
    table[x=step, y=b, y error plus expr=\thisrow{hi}-\thisrow{b}, y error minus expr=\thisrow{b}-\thisrow{lo}]
    {figdata/bachelor-2/30-multiple-regression/out/steps.csv};
\end{axis}
\end{tikzpicture}
\end{openingchart}

% figdata: steps (4.10, 5.22, 3.99, 3.82, 3.71, 3.83, 2.04)
With no control, a year of schooling goes with log earnings 4.10 per cent higher. Adding age raises the estimate to 5.22;
adding age squared brings it to 3.99; race and marriage move it little; adding 1975 earnings halves it, to 2.04. Each
number is a correct answer to a different question. Which question is the right one is not a matter of statistics, and the
last step is a warning taken up in \cref{ch:b2:specification}.

\begin{recall}
From the Year 1 volume: the ordinary least squares estimator with one regressor, fitted values and residuals, the error
term, the zero conditional mean assumption, $R^2$, the omitted variable, estimators, standard errors and confidence
intervals. From mathematics: matrices, the transpose, the inverse, the rank, and orthogonal projection.
\end{recall}

\section{The model in matrix form}\label{sec:b2:multiple-regression:model}

\begin{definition}[Multiple regression model]\label{def:b2:multiple-regression:model}
The \emph{multiple regression model}\index{multiple regression model} writes an outcome as a linear function of several
regressors plus an error: $y_i = \beta_0 + \beta_1x_{i1} + \dots + \beta_{k-1}x_{i,k-1} + u_i$, $i = 1, \dots, n$.
\end{definition}

\begin{definition}[Regressor matrix]\label{def:b2:multiple-regression:regressor-matrix}
The \emph{regressor matrix}\index{regressor matrix} $X$ is the $n \times k$ matrix whose row $i$ is $(1, x_{i1}, \dots,
x_{i,k-1})$: one row per observation, one column per regressor, the first column a column of ones.
\end{definition}

With $y$ and $u$ the $n$-vectors of outcomes and errors and $\beta$ the $k$-vector of coefficients, the model is
$y = X\beta + u$. Three assumptions make $\beta$ the object least squares estimates: the model is linear in $\beta$; $X$
has full column rank $k$, so that no column is a linear combination of the others; and $E[u \mid X] = 0$, so that the
errors have mean zero whatever the regressors. The coefficient $\beta_j$ is then the change in the conditional mean of
$y$ when $x_j$ rises by one and the other regressors are held fixed.

\section{The least-squares estimator}\label{sec:b2:multiple-regression:ols}

Least squares chooses $b$ to minimise the sum of squared residuals $(y - Xb)'(y - Xb) = y'y - 2b'X'y + b'X'Xb$. Setting
the gradient $-2X'y + 2X'Xb$ to zero gives the following equations.

\begin{definition}[Normal equations]\label{def:b2:multiple-regression:normal-equations}
The \emph{normal equations}\index{normal equations} of least squares are $X'Xb = X'y$; when $X$ has full column rank,
$X'X$ is invertible and they have the single solution $\hat\beta = (X'X)^{-1}X'y$.
\end{definition}

The \omterm{def:b2:multiple-regression:normal-equations}{normal equations} say $X'(y - X\hat\beta) = X'e = 0$: the residual vector $e$ is orthogonal to every column of $X$. So
the fitted vector $\hat y = X\hat\beta$ is the orthogonal projection of $y$ on the space spanned by the columns of $X$, and
$e$ is what is left, at right angles to that space. With a column of ones in $X$, the residuals sum to zero.

\begin{omfigure}
\begin{tikzpicture}[econ, every node/.style={font=\small}]
  \fill[omDef!10] (0,0) -- (6,0) -- (7.5,1.8) -- (1.5,1.8) -- cycle;
  \node at (6.2,1.45) {columns of $X$};
  \draw[->, thick] (2,0.6) -- (4.6,1.1) node[midway, below] {$\hat y = X\hat\beta$};
  \draw[->, thick, omDef] (2,0.6) -- (4.6,3.6) node[midway, above left] {$y$};
  \draw[->, thick, omThm] (4.6,1.1) -- (4.6,3.6) node[midway, right] {$e = y - \hat y$};
  \draw (4.6,1.35) -- (4.38,1.31) -- (4.38,1.06);
\end{tikzpicture}

{\small Least squares as projection. The fitted vector is the point of the column space closest to $y$; the residual is
orthogonal to the column space, which is what the \omterm{def:b2:multiple-regression:normal-equations}{normal equations} say.}
\end{omfigure}

The listings use the CPS comparison sample described in the box. The first loads it; the second solves the \omterm{def:b2:multiple-regression:normal-equations}{normal
equations} with numpy and compares the result with statsmodels.

\begin{measure}[The comparison sample]
% ledger: b2:nsw-cps, b1:lalonde1986, b1:dehejia1999; figdata: sample (share capped 0.259)
LaLonde (1986) asked whether econometric methods applied to survey data could reproduce what a field experiment had found
about a training programme. To do so he built comparison samples of men from two American surveys, the Current Population
Survey and the Panel Study of Income Dynamics. Dehejia and Wahba (1999) re-examined the same samples and distribute them;
the CPS one has 15,992 men, with their age, schooling, race, marital status, whether they have a high-school diploma, and
their earnings in 1974, 1975 and 1978. In the file, 25.9 per cent of the men have 1978 earnings equal to the largest value,
25,564.7: earnings above some level were recorded at that level, so the top of the distribution is cut off. This chapter
uses the 13,055 men with positive earnings in both 1975 and 1978. Dropping the men with no earnings in 1978 selects on the
outcome itself, so the coefficients describe men in work, not all men; the experiment itself returns in \cref{ch:b2:causality-and-experiments}.
\end{measure}

\omcode{code/bachelor-2/30-multiple-regression/listings/load_cps.py}{1}{18}{Loading the comparison sample and the
regression sample.}\label{lst:b2:multiple-regression:load}

\omoutput{code/bachelor-2/30-multiple-regression/outputs/load_cps.txt}{A quarter of the men are at the earnings cap; the
regression sample keeps 13,055 men.}\label{out:b2:multiple-regression:load}

\omcode{code/bachelor-2/30-multiple-regression/listings/ols_by_hand.py}{1}{25}{The \omterm{def:b2:multiple-regression:normal-equations}{normal equations} solved with numpy,
against statsmodels.}\label{lst:b2:multiple-regression:ols}

\omoutput{code/bachelor-2/30-multiple-regression/outputs/ols_by_hand.txt}{The same coefficients either way; the residuals
are orthogonal to every column.}\label{out:b2:multiple-regression:ols}

% code: listings/ols_by_hand.py
Holding age, race and marriage fixed, a year of schooling goes with earnings 3.83 per cent higher; married men earn 26.3
log points more than unmarried men of the same age, schooling and race. Age enters with a positive coefficient and a
negative one on its square: earnings rise with age at a falling rate.

\section{Partialling out}\label{sec:b2:multiple-regression:fwl}

What does ``holding the others fixed'' do to the data? Split the regressors into $x$, the one of interest, and $W$, the
others (with the column of ones), so that $y = x\beta_1 + W\gamma + u$.

\begin{definition}[Frisch--Waugh--Lovell theorem]\label{def:b2:multiple-regression:fwl}
The \emph{Frisch--Waugh--Lovell theorem}\index{Frisch--Waugh--Lovell theorem} states that the least-squares coefficient
on $x$ in the regression of $y$ on $x$ and $W$ equals the coefficient in the regression of $\tilde y$ on $\tilde x$, where
$\tilde y$ and $\tilde x$ are the residuals from regressing $y$ and $x$ on $W$.
\end{definition}

\begin{proposition}[Frisch--Waugh--Lovell]\label{prop:b2:multiple-regression:fwl}
Let $M = I - W(W'W)^{-1}W'$, $\tilde x = Mx$ and $\tilde y = My$. Then the coefficient on $x$ in the regression of $y$ on
$(x, W)$ is $\hat\beta_1 = \tilde x'\tilde y/\tilde x'\tilde x = \tilde x'y/\tilde x'\tilde x$.
\end{proposition}

\begin{proof}
$M$ is symmetric, $MM = M$, and $MW = 0$. Write $y = x\hat\beta_1 + W\hat\gamma + e$, with $e$ orthogonal to $x$ and to
every column of $W$. Then $Me = e$, since $e$ is already orthogonal to $W$. Multiply by $M$: $My = Mx\hat\beta_1 + e$.
Multiply by $x'$: $x'My = x'Mx\,\hat\beta_1 + x'e = x'Mx\,\hat\beta_1$. Since $x'My = (Mx)'(My) = \tilde x'\tilde y$ and
$x'Mx = \tilde x'\tilde x$, $\hat\beta_1 = \tilde x'\tilde y/\tilde x'\tilde x$; and $\tilde x'\tilde y = \tilde x'My =
(M\tilde x)'y = \tilde x'y$.
\end{proof}

The coefficient on schooling uses only the variation in schooling that age, race and marriage do not predict. The rest of
the variation is discarded.

\omcode{code/bachelor-2/30-multiple-regression/listings/fwl.py}{1}{29}{Partialling out the controls gives the same
coefficient as the long regression.}\label{lst:b2:multiple-regression:fwl}

\omoutput{code/bachelor-2/30-multiple-regression/outputs/fwl.txt}{Three routes, one coefficient; the controls remove an
eighth of the variance of schooling.}\label{out:b2:multiple-regression:fwl}

\section{The omitted-variable bias formula}\label{sec:b2:multiple-regression:ovb}

\begin{definition}[Omitted-variable bias]\label{def:b2:multiple-regression:ovb}
\emph{Omitted-variable bias}\index{omitted-variable bias} is the difference between the coefficient on a regressor in a
regression that leaves out a variable and the coefficient in the regression that includes it.
\end{definition}

Let the long regression be $y = x\hat\beta_1 + W\hat\beta_W + z\hat\gamma + e$ and the short one leave $z$ out. Let
$\hat\delta$ be the coefficient on $x$ when $z$ itself is regressed on $x$ and $W$. Then, exactly in any sample,
\begin{equation}\label{eq:b2:multiple-regression:ovb}
\hat\beta_1^{\text{short}} = \hat\beta_1^{\text{long}} + \hat\gamma\,\hat\delta.
\end{equation}
Proof: regress both sides of $y = x\hat\beta_1 + W\hat\beta_W + z\hat\gamma + e$ on $x$ and $W$. Least squares is linear
in the outcome, the regression of $x$ and of $W$ on themselves gives coefficient one on themselves, $e$ is orthogonal to
$x$ and $W$ so its coefficients are zero, and $z$ gives $\hat\delta$ on $x$ (Exercise 11).

The formula needs only two signs: the effect of the omitted variable on the outcome and its association with the
regressor.

\begin{omfigure}
\begin{tabular}{lcc}
\toprule
 & $z$ and $x$ positively associated ($\delta > 0$) & negatively associated ($\delta < 0$) \\
\midrule
$z$ raises $y$ ($\gamma > 0$) & short coefficient too high & short coefficient too low \\
$z$ lowers $y$ ($\gamma < 0$) & short coefficient too low & short coefficient too high \\
\bottomrule
\end{tabular}

\smallskip
{\small The sign of the \omterm{def:b2:multiple-regression:ovb}{omitted-variable bias} $\gamma\delta$. Adding age in the opening chart raised the estimate: earnings
rise with age, and older men in the sample have less schooling.}
\end{omfigure}

\omcode{code/bachelor-2/30-multiple-regression/listings/ovb.py}{1}{28}{The \omterm{def:b2:multiple-regression:ovb}{omitted-variable bias} formula checked on the
sample.}\label{lst:b2:multiple-regression:ovb}

\omoutput{code/bachelor-2/30-multiple-regression/outputs/ovb.txt}{Adding 1975 earnings: $0.020429 + 0.399801 \times 0.044806
= 0.038343$.}\label{out:b2:multiple-regression:ovb}

% code: listings/ovb.py
A man with 1 per cent higher earnings in 1975, at the same schooling, age, race and marriage, earns 0.40 per cent more in
1978; and a year of schooling goes with 4.5 per cent higher 1975 earnings. The product, 1.8 points, is the part of the
schooling coefficient that 1975 earnings take over. Whether it should be taken over is another matter: 1975 earnings are
themselves partly a result of schooling, so controlling for them removes part of the effect being measured.

\section{Near-collinear regressors}\label{sec:b2:multiple-regression:collinearity}

\begin{definition}[Perfect collinearity]\label{def:b2:multiple-regression:collinearity}
Regressors are \emph{perfectly collinear}\index{perfect collinearity} when one column of $X$ is an exact linear
combination of the others; $X$ then has rank below $k$, $X'X$ is singular, and the coefficients are not identified.
\end{definition}

A common case is a full set of indicator variables with a constant: an indicator for married men, one for unmarried men,
and the column of ones, which is their sum. Near-collinearity is not a failure of the assumptions, but it costs precision.
Under the assumptions of \cref{ch:b2:regression-inference}, with $R_j^2$ the $R^2$ of the regression of $x_j$ on the
other regressors and $s_j^2$ the sample variance of $x_j$,
\begin{equation}\label{eq:b2:multiple-regression:var}
\operatorname{Var}(\hat\beta_j \mid X) = \frac{\sigma^2}{ns_j^2(1 - R_j^2)}.
\end{equation}
The formula follows from the \omterm{def:b2:multiple-regression:fwl}{Frisch--Waugh--Lovell theorem}: $\hat\beta_j$ is a regression on $\tilde x_j$, whose sum of
squares is $ns_j^2(1 - R_j^2)$ (Exercise 12).

\begin{definition}[Variance inflation factor]\label{def:b2:multiple-regression:vif}
The \emph{variance inflation factor}\index{variance inflation factor} of regressor $j$ is $1/(1 - R_j^2)$: the factor by
which the variance of $\hat\beta_j$ exceeds what it would be if $x_j$ were uncorrelated with the other regressors.
\end{definition}

\omcode{code/bachelor-2/30-multiple-regression/listings/collinearity.py}{1}{20}{Simulated regressors with rising
correlation, and an exactly collinear design.}\label{lst:b2:multiple-regression:collinearity}

\omoutput{code/bachelor-2/30-multiple-regression/outputs/collinearity.txt}{The standard error grows with the correlation
roughly as the square root of the VIF; with an exact multiple the matrix loses a rank.}\label{out:b2:multiple-regression:collinearity}

% code: listings/collinearity.py
At a correlation of 0.99 the \omterm{def:b2:multiple-regression:vif}{variance inflation factor} is 48.5 and the standard error of $\hat\beta_1$ is 0.311, about
seven times the 0.045 of uncorrelated regressors; the estimate itself is still unbiased. Near-collinearity is a shortage
of information about the separate effects, not an error to be corrected; the remedies are more data or a question that
does not need the separate effects.

\section{Fit}\label{sec:b2:multiple-regression:fit}

The $R^2$ of a multiple regression is $1 - e'e/\sum_i(y_i - \bar y)^2$, as with one regressor. Adding a regressor never
lowers it, since least squares could always set the new coefficient to zero.

\begin{definition}[Adjusted R-squared]\label{def:b2:multiple-regression:adjusted-r2}
The \emph{adjusted $R^2$}\index{adjusted R-squared} is $\bar R^2 = 1 - \frac{e'e/(n - k)}{\sum_i(y_i - \bar y)^2/(n -
1)} = 1 - (1 - R^2)\frac{n - 1}{n - k}$; it can fall when a regressor is added.
\end{definition}

% code: listings/ols_by_hand.py; figdata: steps
In the listing $R^2$ is 0.1181 and $\bar R^2$ 0.1177: with 13,055 observations and seven coefficients the correction is
tiny. Neither number says whether the coefficient on schooling is the effect of schooling. Adding 1975 earnings raises
$R^2$ from 0.118 to 0.254 and moves the schooling coefficient from 3.83 to 2.04; a better fit is not a better answer to a causal
question.

\begin{model}{The linear regression model}
\assumes The conditional mean of the outcome is linear in the regressors; the \omterm{def:b2:multiple-regression:regressor-matrix}{regressor matrix} has full column rank; the
error has mean zero given the regressors.
\predicts Least squares estimates each coefficient without bias: the change in the mean outcome when one regressor rises
by one and the others are held fixed.
\tested Linearity can be checked against more flexible forms (\cref{ch:b2:specification}); the zero conditional mean
assumption cannot be tested from the residuals, which are orthogonal to the regressors by construction, but can be tested
against an experiment, as LaLonde did (\cref{ch:b2:causality-and-experiments}).
\breaks When a variable that affects the outcome and is correlated with a regressor is left out, the coefficient absorbs
its effect, by the amount $\gamma\delta$; and when a control is itself an outcome of the regressor, holding it fixed
removes part of the effect.
\end{model}

\section{Exercises}\label{sec:b2:multiple-regression:exercises}

\begin{exercise}[$\star$]\label{exo:b2:multiple-regression:1}
Read the opening chart. Give the estimate with no control, after adding age, and with every control. Which single step
moves the estimate most?
\end{exercise}

\begin{exercise}[$\star$]\label{exo:b2:multiple-regression:2}
A quarter of the men have 1978 earnings at the cap. What does the cap do to the measured differences in log earnings
between men with much schooling and men with little?
\end{exercise}

\begin{exercise}[$\star$]\label{exo:b2:multiple-regression:3}
Compute $\bar R^2$ from $R^2 = 0.1181$, $n = 13{,}055$ and $k = 7$.
\end{exercise}

\begin{exercise}[$\star$]\label{exo:b2:multiple-regression:4}
1975 earnings are measured after most of the men had finished school. Why might controlling for them be a mistake? What
would show that a control is a bad one?
\end{exercise}

\begin{exercise}[$\star\star$]\label{exo:b2:multiple-regression:5}
By hand: with five observations, $x = (1, 2, 3, 4, 5)$, $d = (0, 0, 1, 1, 1)$ and $y = (4, 4, 9, 13, 14)$, regress $y$ on
a constant, $x$ and $d$. Write $X'X$ and $X'y$, solve the \omterm{def:b2:multiple-regression:normal-equations}{normal equations}, and compute $R^2$ and $\bar R^2$.
\end{exercise}

\begin{exercise}[$\star\star$]\label{exo:b2:multiple-regression:6}
Predict, then compute: adding age raises the schooling coefficient from 4.10 to 5.22. The correlation between age and
schooling in the sample is $-0.15$; age has a coefficient of 0.0192 in the long regression, and the regression of age on
schooling has slope $-0.585$. Check the \omterm{def:b2:multiple-regression:ovb}{omitted-variable bias} formula.
\end{exercise}

\begin{exercise}[$\star\star$]\label{exo:b2:multiple-regression:7}
\textbf{(D)} (a) With the Brenland household survey, regress the log gross income of heads with positive income on years
of schooling, then add age and age squared. (b) With the CPS sample, add \texttt{nodegree} to the regression of the
listing \texttt{ols\_by\_hand.py}. What happens to the schooling coefficient, its standard error, and its \omterm{def:b2:multiple-regression:vif}{variance
inflation factor}?
\end{exercise}

\begin{exercise}[$\star\star$]\label{exo:b2:multiple-regression:8}
Change $n$ to 2000 in the collinearity listing. By what factor should the standard errors fall? Compare with the output.
\end{exercise}

\begin{exercise}[$\star\star$]\label{exo:b2:multiple-regression:9}
A regression includes a constant, an indicator for married men and an indicator for unmarried men. Show that $X$ does not
have full column rank. Give two ways to fix it, and say what the coefficients then mean.
\end{exercise}

\begin{exercise}[$\star\star\star$]\label{exo:b2:multiple-regression:10}
Show that $M = I - W(W'W)^{-1}W'$ is symmetric and idempotent, that $MW = 0$, and that $My$ is the residual vector of the
regression of $y$ on $W$.
\end{exercise}

\begin{exercise}[$\star\star\star$]\label{exo:b2:multiple-regression:11}
Prove \eqref{eq:b2:multiple-regression:ovb} in full.
\end{exercise}

\begin{exercise}[$\star\star\star$]\label{exo:b2:multiple-regression:12}
Using the \omterm{def:b2:multiple-regression:fwl}{Frisch--Waugh--Lovell theorem} and $\operatorname{Var}(u \mid X) = \sigma^2I$, derive
\eqref{eq:b2:multiple-regression:var}.
\end{exercise}

\section{Problem: Brenland's household survey}\label{sec:b2:multiple-regression:problem}

\begin{problem}[{Weekend problem --- Brenland's household survey: matrices, the return to schooling with controls,
partialling out and omitted-variable bias, and the collinearity of age and experience}]
\label{pb:b2:multiple-regression:1}
Brenland's household survey of 2025 has 2,000 households; 1,922 heads have positive gross income. Let $y_i$ be the log
gross income of head $i$, $s_i$ his years of schooling and $a_i$ his age.

\medskip
\noindent\textbf{Part I --- Matrices.}
\begin{enumerate}
  \item Write the \omterm{def:b2:multiple-regression:regressor-matrix}{regressor matrix} of the regression of $y$ on a constant, $s$, $a$ and $a^2$: its size and a typical row.
  \item What is the size of $X'X$? Of $X'y$?
  \item What does the first element of $X'X$ count? The first element of $X'y$?
  \item Why must $X'X$ be invertible for the estimator to exist?
  \item What does the first normal equation say about the residuals?
\end{enumerate}

\medskip
\noindent\textbf{Part II --- Returns to schooling with controls.}
\begin{enumerate}[resume]
  \item Regress $y$ on a constant and $s$. Give the coefficient on $s$ and $R^2$.
  \item Add $a$ and $a^2$. Give the coefficient on $s$, its 95 per cent confidence interval, $R^2$ and $\bar R^2$.
  \item Interpret the coefficient on $s$ in the long regression.
  \item At what age is predicted income highest, holding schooling fixed?
  \item Why does the coefficient on $s$ barely move?
\end{enumerate}

\medskip
\noindent\textbf{Part III --- Partialling out and omitted-variable bias.}
\begin{enumerate}[resume]
  \item Regress $s$ on a constant, $a$ and $a^2$, and keep the residual $\tilde s$. Compare the variances of $s$ and $\tilde s$.
  \item Regress $y$ on $\tilde s$. What coefficient do you get?
  \item Give $\delta_a$ and $\delta_{a^2}$, the coefficients on $s$ when $a$ and $a^2$ are regressed on a constant and $s$.
  \item Write the short coefficient as the long one plus the bias, with two omitted variables.
  \item Compute the bias and check the identity.
  \item Why are $\delta_a$ and $\delta_{a^2}$ so different in size, and yet give a bias near zero?
\end{enumerate}

\medskip
\noindent\textbf{Part IV --- Collinearity of age and experience.}
\begin{enumerate}[resume]
  \item Define experience as $a - s - 6$. Regress $y$ on a constant, $s$, $a$ and experience. What goes wrong?
  \item Give the rank of the \omterm{def:b2:multiple-regression:regressor-matrix}{regressor matrix}.
  \item Which combinations of the coefficients can still be estimated?
  \item Give the \omterm{def:b2:multiple-regression:vif}{variance inflation factors} of $s$, $a$ and $a^2$ in the regression of Part II.
  \item Why is the high VIF of $a$ and $a^2$ no problem for the coefficient on $s$?
  \item With half the sample, by what factor would the standard error of the schooling coefficient grow?
  \item If ability raised income and was positively correlated with schooling, which way would the estimate be biased?
  \item Could controlling for hours worked fix that bias? Could it create another?
  \item State the result: the return to schooling with its interval, and the size of the \omterm{def:b2:multiple-regression:ovb}{omitted-variable bias} from age.
\end{enumerate}
\end{problem}
